\section{Appendix}

\subsection{Benchmark Splits and Evaluation Protocol}
The benchmark has one ID training split, sand, and four zero-shot evaluation splits: Sand-Hard, Sand-Soft, cohesive snow, and elastoplastic. Main results also report ID-sand evaluation, so the tables cover five evaluation splits in total. All methods train on sand only, and OOD scores therefore measure extrapolation under hidden material dynamics rather than interpolation over a labeled material mixture. M1 and M7 main-result means use three training seeds and ten deterministic evaluation episodes per seed. M8 is reported as a single-seed diagnostic privileged-parameter baseline. The reported metrics are transfer efficiency, spill ratio, and success rate, where success means exceeding 30\% transfer.

\subsection{Probe-Then-Act Architecture Details}
\PTA{} uses a two-stage episode with a fixed open-loop probing segment of length $H_p=3$ followed by a residual PPO task policy. During probing, the task residual is held at zero and the scripted probe trace is recorded. A two-layer MLP encoder with hidden size 256 maps the trace to a 16-dimensional latent mean and learned scale; the policy conditions on the current state together with this latent pair. We call this pair a belief for brevity, but it is task-conditioning information rather than a calibrated Bayesian posterior. The probe is fixed in this paper and is not optimized by an information-gain objective.

\subsection{Implementation and Environment Details}
All compared systems use 500k-step PPO training budgets with matched hyperparameters. The policy/value networks use MLP heads of size $[256,256]$, learning rate $3\times 10^{-4}$, batch size 64, 10 epochs per update, and $\gamma=0.99$. The Genesis MPM benchmark uses grid density 128 with 25 substeps per control step and control $\Delta t = 2\times10^{-3}$~s. The action is a 7-DoF joint-space residual with residual scale 0.05 added to a 410-step scripted edge-push trajectory followed by an 80-step settle.

\subsection{Ablation Interpretation}
The no-probe and no-belief ablations are interpreted relative to the elastoplastic transfer gain. Removing the probe drops transfer from 60.7\% to 33.1\%, while removing the encoder output drops transfer to 46.3\%, nearly eliminating the full \PTA{} transfer gain over the 46.0\% reactive baseline. The auxiliary spill and success columns should be read descriptively: M7 Full has the lowest spill and highest success in this table, but No-Belief is close to M1 on transfer and does not dominate M1 on every metric. These ablations support the component interpretation for the elastoplastic transfer gain in this architecture and split; they do not establish that active probing is generally useful across all material families.

\subsection{Privileged-Parameter Baseline}
The privileged-parameter baseline receives ground-truth material parameters concatenated to the policy observation but is trained in the same sand-only residual-control setting as the other compared systems. Its table row is a single-seed diagnostic observation: M8 is strong on ID sand, snow, and Sand-Hard, trails M1 on Sand-Soft, and scores 0.0\% transfer on elastoplastic. In this setup, passive parameter access did not produce a useful elastoplastic response curve, whereas an in-episode probe trace did. This is not a universal claim that interaction traces dominate privileged information.

\subsection{Failure Modes and Scope}
The main failure mode is probe-induced degradation on irreversible or partly irreversible substrates. On ID sand, Sand-Hard, Sand-Soft, and snow, \PTA{} underperforms the reactive baseline; the paper therefore makes a narrow positive claim only on the elastoplastic split and explicitly avoids a broad robustness claim. The recoverability account is an explanatory hypothesis for this sign pattern, not a validated law. The current benchmark does not test additional recoverable material families, learned probe policies, probe-gating mechanisms, realistic tactile sensors, or sim-to-real deployment.
